\section{Related Work}\label{sec:related}

\subsection{Search-Based and Random Test Generation}

Automated unit-test generation predates language models by more than a decade. Randoop
builds test sequences incrementally and uses feedback from executing each prefix to discard
sequences that are illegal or redundant~\cite{pacheco2007randoop}. EvoSuite treats an entire
test suite as one individual in a genetic search and evolves it against a coverage-based
fitness function~\cite{fraser2011evosuite}. Both have been evaluated at scale: EvoSuite
reached high branch coverage across a large sample of open-source classes, while also
showing that a substantial fraction of classes remain hard for reasons unrelated to the
search itself~\cite{fraser2014largescale}. Pynguin ports the search-based approach to
Python, where the absence of static types removes much of the information the Java
generators depend on~\cite{lukasczyk2022pynguin,lukasczyk2023empirical}.

Two findings from this line of work bear directly on our design. First, coverage and
fault detection are not interchangeable: Shamshiri et al.\ found that automatically
generated suites with high coverage detected only a minority of real
faults~\cite{shamshiri2015realfaults}, and Inozemtseva and Holmes report that coverage
correlates with suite effectiveness far more weakly than is commonly
assumed~\cite{inozemtseva2014coverage}. Second, the difficulty of a target is not captured
by its complexity alone --- constructing the receiver object is frequently the binding
constraint. We return to both points in Sections~\ref{sec:tiers} and~\ref{sec:discussion}.

\subsection{LLM-Based Test Generation}

Code-trained language models made test generation a prompting problem~\cite{chen2021codex},
and the earliest work in this direction fine-tuned transformers on method--test pairs with
surrounding focal context~\cite{tufano2020athenatest}. Subsequent systems moved to
prompting. TestPilot generates tests for JavaScript functions from documentation and usage
examples, and re-prompts with the failure message when a test does not
pass~\cite{schafer2024testpilot}. ChatTester adds an iterative repair loop driven by
compiler and runtime errors, reporting large gains over a single-shot
baseline~\cite{yuan2023chattester}. Siddiq et al.\ evaluate several models on JUnit
generation and find compilation failure --- not weak assertions --- to be the dominant
failure mode~\cite{siddiq2024junit}. A recent survey catalogues the space and notes that
evaluation practice is uneven across the field~\cite{wang2024llmtesting}.

Our study is deliberately narrower than these systems in one respect and stricter in
another. We use a \emph{single} API call per function, with no repair loop and no
retrieval, because a repair loop conflates the quality of the initial generation with the
quality of the error signal fed back into it; when compilation failure is itself the object
of study, that conflation is fatal. Where those systems ask how high the success rate can
be driven, we ask what the initial generation is actually missing, and whether it is
something a prompt can supply. The two prompt conditions of
Section~\ref{sec:generation} differ in exactly one variable so that the answer is
attributable to that variable.

\subsection{Evaluation Metrics and Their Validity}

Mutation analysis is the standard fault-based criterion for test-suite adequacy, and
mutants have been shown to be a reasonable --- if imperfect --- substitute for real faults
in controlled comparisons~\cite{just2014mutants,papadakis2019mutation}. It is
comparatively expensive, which is why much of the LLM test-generation literature reports
compilation success and line or branch coverage instead.

This substitution is where our contribution sits. Compilation success and coverage are
cheap because they ask less, and what they ask can be satisfied without testing anything: a
reflective call compiles regardless of whether the named method exists, and a suite that
constructs no object and calls nothing still passes. Neither weakness is hypothetical; both
occur in the corpus measured here (Section~\ref{sec:tiers}). We therefore report four
nested tiers and rest every claim on the only tier that cannot be satisfied vacuously.

A second and less discussed threat concerns the benchmark rather than the metric. Sampling
targets by cyclomatic complexity~\cite{mccabe1976complexity} is standard practice, but
complexity does not imply reachability: a method may be private, may be one of several
overloads sharing a name, or may belong to a class that cannot be instantiated from
outside. A suite that fails on such a target tells us about the sampling, not about the
generator. To our knowledge, the fraction of a complexity-sampled benchmark that is
actually addressable from an external test is rarely reported. We report the full sampling
funnel for this reason, and the attrition turns out to be the study's most transferable
number.

\subsection{Methodological Practice}

Because the present work re-examines a dataset after having seen results obtained on an
earlier one, the distinction between repairing an instrument and selecting a favourable
outcome is central. Kerr's account of hypothesising after the results are
known~\cite{kerr1998harking} describes the failure mode precisely. We address it in two
ways: the sampling criteria, tool versions, budgets and hypotheses were committed to the
repository before any measurement was run, and the earlier study's dataset and reported
numbers are left unchanged --- this paper stands beside it rather than replacing it. For
the statistical comparisons we follow the guidance of Arcuri and Briand on non-parametric
testing and effect sizes for randomised algorithms in software
engineering~\cite{arcuri2014hitchhiker}.
